% chapter4.tex
\chapter{روش‌شناسی پژوهش و تصریح مدل}

\section{مقدمه}
(در مقدمه این فصل توضیحات کلی در مورد محتوای فصل ارائه می‌شود.)

\section{طرح، روش و یا رویکرد پژوهش}
(در این بخش با توجه به اهداف پژوهش، طرح، رویکرد، و یا روش پژوهش و نوع مدل اقتصادسنجی مشخص می‌شود.)

\section{جامعه پژوهش و نمونه‌گیری}
(در این بخش جامعه آماری پژوهش توصیف می‌شود. منظور از جامعه پژوهش مجموعه اعضای حقیقی یا فرضی مورد مطالعه است.)

\section{روش تجزیه و تحلیل داده‌ها و تصریح مدل}
(تجزیه و تحلیل داده‌ها فرآیندی است که به استخراج اطلاعات ارزشمند از داده‌ها اشاره دارد. فرمول‌های اقتصادسنجی و آزمون‌های تشخیصی در این بخش قرار می‌گیرند.)